\documentclass[11pt,letterpaper]{article}

% ── Geometry ──────────────────────────────────────────────────────────────────
\usepackage[margin=1in]{geometry}

% ── Pagination ────────────────────────────────────────────────────────────────
\widowpenalty=10000          % No widow lines (last line alone at page top)
\clubpenalty=10000           % No orphan lines (first line alone at page bottom)
\brokenpenalty=10000         % No hyphenated words across page breaks
\raggedbottom                % Allow uneven page bottoms rather than stretching

% ── Fonts ─────────────────────────────────────────────────────────────────────
\usepackage[T1]{fontenc}
\usepackage{newpxtext}     % Palatino-based serif (TeX Gyre Pagella)
\usepackage{newpxmath}     % Matching math font

% ── Colors ────────────────────────────────────────────────────────────────────
\usepackage[dvipsnames]{xcolor}
\definecolor{SectionBlue}{HTML}{1B3A5C}
\definecolor{AccentGray}{HTML}{4A4A4A}
\definecolor{QuoteBorder}{HTML}{8B9DAF}
\definecolor{RiskRed}{HTML}{C0392B}
\definecolor{RiskAmber}{HTML}{E67E22}
\definecolor{RiskGreen}{HTML}{27AE60}
\definecolor{BoxBg}{HTML}{F7F8FA}

% ── Tables ────────────────────────────────────────────────────────────────────
\usepackage{booktabs}
\usepackage{tabularx}

% ── Boxes ─────────────────────────────────────────────────────────────────────
\usepackage[framemethod=tikz]{mdframed}

% ── Lists ─────────────────────────────────────────────────────────────────────
\usepackage{enumitem}

% ── Hyperlinks ────────────────────────────────────────────────────────────────
\usepackage{hyperref}
\hypersetup{
  colorlinks=true,
  linkcolor=SectionBlue,
  urlcolor=SectionBlue,
  citecolor=SectionBlue,
  pdfauthor={},
  pdftitle={PR/FAQ},
  pdfsubject={Working Backwards — Product Discovery},
}

% ── Bibliography ─────────────────────────────────────────────────────────────
\usepackage[style=numeric,backend=biber,sorting=none]{biblatex}
% \addbibresource{prfaq.bib}  % Set filename when using citations

% ── Section styling ───────────────────────────────────────────────────────────
\usepackage{titlesec}
\titleformat{\section}
  {\Large\bfseries\color{SectionBlue}}
  {}
  {0pt}
  {}
  [\vspace{-0.5em}\textcolor{QuoteBorder}{\rule{\textwidth}{0.4pt}}]

\titleformat{\subsection}
  {\large\bfseries\color{AccentGray}}
  {}
  {0pt}
  {}

% ══════════════════════════════════════════════════════════════════════════════
% Custom environments
% ══════════════════════════════════════════════════════════════════════════════

% ── \prsection{title} — styled press release section header ──────────────────
\newcommand{\prsection}[1]{%
  \vspace{1em}%
  {\large\bfseries\color{SectionBlue}#1}%
  \vspace{0.3em}\par%
}

% ── customerquote — indented, italic customer quote ──────────────────────────
\newmdenv[
  topline=false,
  bottomline=false,
  rightline=false,
  linewidth=3pt,
  linecolor=QuoteBorder,
  backgroundcolor=BoxBg,
  innerleftmargin=12pt,
  innerrightmargin=12pt,
  innertopmargin=8pt,
  innerbottommargin=8pt,
  skipabove=\baselineskip,
  skipbelow=\baselineskip,
]{customerquote}

% ── spokespersonquote — indented spokesperson quote ──────────────────────────
\newmdenv[
  topline=false,
  bottomline=false,
  rightline=false,
  linewidth=3pt,
  linecolor=SectionBlue,
  backgroundcolor=BoxBg,
  innerleftmargin=12pt,
  innerrightmargin=12pt,
  innertopmargin=8pt,
  innerbottommargin=8pt,
  skipabove=\baselineskip,
  skipbelow=\baselineskip,
]{spokespersonquote}

% ── faqpair — Q&A pair with bold question ────────────────────────────────────
% Label FAQ pairs that explain judgment calls: \begin{faqpair}{Question}\label{faq:slug}
% Then cross-reference from the press release: (see p.\,\pageref{faq:slug})
\newenvironment{faqpair}[1]{%
  \vspace{0.5em}%
  \noindent\textbf{\color{SectionBlue}Q: #1}\par%
  \vspace{0.2em}%
  \begin{adjustwidth}{1em}{0pt}%
  \setlength{\parindent}{0pt}%
}{%
  \end{adjustwidth}%
  \vspace{0.5em}%
}

% ── adjustwidth (for faqpair indentation) ────────────────────────────────────
\usepackage{changepage}

% ══════════════════════════════════════════════════════════════════════════════
% Document
% ══════════════════════════════════════════════════════════════════════════════

\begin{document}

% ── Title block ───────────────────────────────────────────────────────────────
\begin{center}
  {\LARGE\bfseries\color{SectionBlue} PRODUCT NAME} \\[0.5em]
  {\large\color{AccentGray} Subtitle — one-line value proposition} \\[1.5em]
  {\small\color{AccentGray} Date \quad|\quad Author \quad|\quad Team}
\end{center}

\vspace{1em}
\textcolor{QuoteBorder}{\rule{\textwidth}{0.8pt}}
\vspace{1em}

% ══════════════════════════════════════════════════════════════════════════════
% PRESS RELEASE
% ══════════════════════════════════════════════════════════════════════════════

\section*{Press Release}

\prsection{Summary}

CITY, STATE — Today, COMPANY announced PRODUCT, which provides BENEFIT
to TARGET CUSTOMER. Unlike ALTERNATIVES, PRODUCT offers DIFFERENTIATOR,
enabling customers to OUTCOME.

\prsection{Problem}

TARGET CUSTOMER struggles with PROBLEM. Today they cope by CURRENT WORKAROUND,
which results in NEGATIVE CONSEQUENCE. Despite EXISTING SOLUTIONS, customers
still face UNRESOLVED PAIN POINT because REASON.

\prsection{Solution}

PRODUCT solves this by APPROACH. Customers can now KEY ACTION, which means
they experience BENEFIT. The product works by BRIEF MECHANISM, requiring
MINIMAL EFFORT from the customer.

\prsection{Customer Quote}

\begin{customerquote}
\textit{``Before PRODUCT, I spent PAIN POINT doing WORKAROUND. Now I can
DESIRED OUTCOME in TIMEFRAME. It has changed how I BROADER IMPACT.''}

\raggedleft --- \textbf{Customer Name}, Role at Company
\end{customerquote}

\prsection{Getting Started}

Getting started with PRODUCT is simple. Customers FIRST STEP, then
SECOND STEP, and within TIMEFRAME they can INITIAL VALUE. There is no
need to REMOVED FRICTION.

\prsection{Spokesperson Quote}

\begin{spokespersonquote}
``We built PRODUCT because we saw that TARGET CUSTOMER deserved
BETTER EXPERIENCE. By focusing on KEY DESIGN DECISION, we've been able to
deliver QUANTIFIABLE IMPROVEMENT while keeping IMPORTANT QUALITY.''

\raggedleft --- \textbf{Spokesperson Name}, Title at Company
\end{spokespersonquote}

\prsection{Call to Action}

To learn more and get started, visit URL. PRODUCT is available starting
DATE with PRICING MODEL.

\newpage

% ══════════════════════════════════════════════════════════════════════════════
% FREQUENTLY ASKED QUESTIONS
% ══════════════════════════════════════════════════════════════════════════════

\section*{Frequently Asked Questions}

% ── External FAQs (Customer-facing) ──────────────────────────────────────────

\subsection*{External FAQs}

\begin{faqpair}{What is PRODUCT and who is it for?}
  PRODUCT is BRIEF DESCRIPTION designed for TARGET CUSTOMER who needs
  to CORE JOB TO BE DONE.
\end{faqpair}

\begin{faqpair}{How is this different from COMPETITOR/ALTERNATIVE?}
  Unlike ALTERNATIVE, PRODUCT provides DIFFERENTIATOR because REASON.
  This means customers can UNIQUE BENEFIT.
\end{faqpair}

\begin{faqpair}{How do I get started?}
  ONBOARDING STEPS. Most customers see INITIAL VALUE within TIMEFRAME.
\end{faqpair}

% ── Internal FAQs (Business-facing) ──────────────────────────────────────────

\subsection*{Internal FAQs}

\subsubsection*{Value \& Market}

\begin{faqpair}{What is the total addressable market?}
  TAM ANALYSIS with supporting data.
\end{faqpair}

\begin{faqpair}{What evidence do we have that customers want this?}
  CUSTOMER EVIDENCE: interviews, surveys, usage data, or market signals.
\end{faqpair}

\subsubsection*{Technical}

\begin{faqpair}{What are the major technical risks?}
  TECHNICAL RISKS and mitigation strategies.
\end{faqpair}

\begin{faqpair}{What is the estimated development timeline?}
  TIMELINE with key milestones and dependencies.
\end{faqpair}

\subsubsection*{Business}

\begin{faqpair}{What is the revenue model?}
  REVENUE MODEL description with pricing rationale.
\end{faqpair}

\begin{faqpair}{What are the key metrics for success?}
  SUCCESS METRICS with targets and measurement approach.
\end{faqpair}

% ══════════════════════════════════════════════════════════════════════════════
% FOUR RISKS ASSESSMENT
% ══════════════════════════════════════════════════════════════════════════════

\section*{Risk Assessment}

\renewcommand{\arraystretch}{1.6}
\begin{tabularx}{\textwidth}{@{} l l X @{}}
\toprule
\textbf{\color{SectionBlue}Risk} & \textbf{\color{SectionBlue}Rating} & \textbf{\color{SectionBlue}Assessment} \\
\midrule
Value      & \textcolor{RiskAmber}{RATING} & Will customers buy or choose to use this? \\
Usability  & \textcolor{RiskGreen}{RATING} & Can customers figure out how to use it? \\
Feasibility & \textcolor{RiskGreen}{RATING} & Can we build it with available resources? \\
Viability  & \textcolor{RiskGreen}{RATING} & Does it work for the business? \\
\bottomrule
\end{tabularx}
\renewcommand{\arraystretch}{1.0}

% ══════════════════════════════════════════════════════════════════════════════
% FEATURE APPENDIX
% ══════════════════════════════════════════════════════════════════════════════

\section*{Feature Appendix}

Defines the scope boundary for the product. Every feature is explicitly categorized to prevent scope creep and force prioritization decisions upfront.

\subsection*{Must Do}

Features that are essential for launch. Without these, the product does not solve the core problem.

\begin{itemize}[leftmargin=1.5em]
  \item FEATURE — RATIONALE
  \item FEATURE — RATIONALE
  \item FEATURE — RATIONALE
\end{itemize}

\subsection*{Should Do}

Features that meaningfully improve the product but are not launch-blocking. Candidates for fast follow-up.

\begin{itemize}[leftmargin=1.5em]
  \item FEATURE — RATIONALE
  \item FEATURE — RATIONALE
\end{itemize}

\subsection*{Won't Do}

Features explicitly excluded. Naming what you won't build is as important as naming what you will — it prevents scope creep and clarifies the product's identity.

\begin{itemize}[leftmargin=1.5em]
  \item FEATURE — WHY NOT
  \item FEATURE — WHY NOT
\end{itemize}

% ══════════════════════════════════════════════════════════════════════════════
% REFERENCES
% ══════════════════════════════════════════════════════════════════════════════

\printbibliography[title={References}]

\end{document}
